\begin{exercice}\label{exo4-10} (\coolexo)

Une variable al\'eatoire $X : \Omega \rightarrow \mathbb{R}$ suit la {\it loi exponentielle} $\mathcal{E} (\lambda)$ de param\`etre $\lambda > 0$ si sa densit\'e est donn\'ee par:

\[
f(x) = \lambda e^{-\lambda x} \mathds{1}_{\mathbb{R}_+} (x),
\]

% \mathbb{\1} surprenant !!

pour $x \in \mathbb{R}$. 

\begin{quotation}
$\rightarrow $ {\it C'est le mod\`ele le plus simple pour la dur\'ee de vie d'un mat\'eriel (sans usure).}
\end{quotation}

\begin{enumerate}

\item V\'erifier que $f$ est bien une densit\'e de probabilit\'e.

\item Calculer la fonction de r\'epartition $F_X (t)$ de $X$ et tracer la courbe correspondante.

\item On appelle transform\'ee de Laplace de la variable la fonction :

\[
\mathcal{L} (s) = E [ e^{-sX} ] = \int_0^{+ \infty} e^{-sx} f(x) dx.
\]

Calculer la transform\'ee de Laplace pour $X$, qui suit la loi exponentielle.

\item Montrer que pour toute variable al\'eatoire $X$, on a :

\[
\mathcal{L}' (0) = - E[X], 
\]

et

\[
\mathcal{L}'' (0) = E[X^2]. 
\]

En d\'eduire la valeur de la moyenne et de la variance pour une variable $X$ qui suit la loi
exponentielle.

\item Montrer que $X$ est effectivement sans m\'emoire (ou sans usure), c'est \`a dire :

\[
P ( X \geq s+t | X \geq t) = P ( X \geq s),
\]

pour $s,t \geq 0$.

\item Calculer $P(X \in (t , t + \delta t) | X \geq t )$ pour $\delta t > 0$ petit. Que conclut-t'on ?

\end{enumerate}

%\corrref{TD1_1}

\end{exercice}
